\chapter{Discussion and conclusions}
\label{ch:discuss}

\section{Limits of validity}
\begin{itemize}[nosep]
\item \textbf{Linear elastic, small strain} everywhere: no plasticity, large rotation, contact or dynamics. Rubber in the 3D tab can exceed $100\%$ strain, where results are indicative only.
\item \textbf{Half-space} (3D field, sections): an infinite body, so a finite body's sides and base are not traction-free; the FEM block quantifies the error (Section~\ref{sec:convergence}).
\item \textbf{Finite elements}: plane stress only; CST and Q4 are stiff in bending until refined; bricks lock for nearly incompressible materials; the dense 2D solver is limited to about 3600 degrees of freedom.
\item \textbf{Environment}: Three.js and Plotly load from CDNs, so the interface needs network access.
\end{itemize}

\section{Conclusions}
\begin{enumerate}[nosep]
\item The project combines half-space elasticity, closed-form simulators and finite elements in one interface, so a result obtained one way can be checked another.
\item All nine simulators reproduce their closed-form solutions to the printed precision (Table~\ref{tab:defaults}) and pass the recorded limiting-case and invariance tests.
\item The FEM solver converges monotonically: to $0.01\%$ for the bar with 32 elements, within $0.3\%$ of beam theory for the Q4 cantilever, and $K_t=3.06$ for a hole in a wide plate against the classical $3$.
\item The half-space approximation of the 3D field is good to about $3\%$ for bodies several times larger than the loaded area.
\end{enumerate}

\section*{Acknowledgements}
\addcontentsline{toc}{section}{Acknowledgements}
We are sincerely grateful to \textbf{Prof. Prakhar Gupta} for giving us the opportunity to work on this project. His ideas shaped the project from the start: the discussions with him while we were deciding what to build gave us its direction and much of its scope. We thank him for his guidance, his encouragement and his trust in letting us develop the idea in our own way.

\begin{flushright}
\sffamily Abhinav Kumar Sudhir Choudhary, Rudra Thummar,\\ Tilak Asodariya and Sriramnarayanan Mohan
\end{flushright}

\begin{thebibliography}{99}\small
\bibitem{timoshenko} S.~P. Timoshenko and J.~N. Goodier, \emph{Theory of Elasticity}, 3rd ed., McGraw-Hill, 1970.
\bibitem{johnson} K.~L. Johnson, \emph{Contact Mechanics}, Cambridge University Press, 1985.
\bibitem{boussinesq} J.~Boussinesq, \emph{Application des potentiels \`a l'\'etude de l'\'equilibre et du mouvement des solides \'elastiques}, Gauthier-Villars, 1885.
\bibitem{cerruti} V.~Cerruti, ``Ricerche intorno all'equilibrio dei corpi elastici isotropi,'' \emph{Atti Accad. Lincei}, 1882.
\bibitem{hibbeler} R.~C. Hibbeler, \emph{Mechanics of Materials}, 10th ed., Pearson, 2017.
\bibitem{cook} R.~D. Cook, D.~S. Malkus, M.~E. Plesha and R.~J. Witt, \emph{Concepts and Applications of Finite Element Analysis}, 4th ed., Wiley, 2002.
\bibitem{newmark} N.~M. Newmark, ``Simplified computation of vertical pressures in elastic foundations,'' \emph{Univ. Illinois Eng. Exp. Station Circular} 24, 1935.
\bibitem{kirsch} G.~Kirsch, ``Die Theorie der Elastizit\"at und die Bed\"urfnisse der Festigkeitslehre,'' \emph{Z. VDI} 42, 1898.
\end{thebibliography}
